For a complete and connected Riemannian manifold $M$, we view its Laplacian $\Delta$
as an unbounded and symmetric operator with domain $C^\infty_{\rm c}(M) \subseteq L^2(M)$.
The closure of $\Delta$, which we also denote by $\Delta$, is a self-adjoint operator in $L^2(M)$.
Its spectrum as a subset of $\R$ will be denoted by $\sigma(M)$.

Functional analysis of self-adjoint operators yields the orthogonal decomposition
\begin{align*}%\label{deco}
L^2(M) = H_{\p}(M) \oplus H_{\ac}(M) \oplus H_{\sc}(M)
\end{align*}
of $L^2(M)$ into $\Delta$-invariant subspaces such that
\begin{enumerate}\itemsep=0pt
\item[(1)] $H_{\p}(M)$ is spanned by eigenfunctions;
\item[(2)] $H_{\ac}(M)$ consists of functions with absolutely continuous spectral measure;
\item[(3)] $H_{\sc}(M)$ consists of functions with singularly continuous spectral measure.
\end{enumerate}
We present some more details about these notions and decompositions in Appendix~\ref{secspec},
since they are needed in our discussion of spectra of Riemannian products.

There is the well known decomposition $\sigma(M)=\sigma_{\d}(M)\cup\sigma_{\ess}(M)$ of the spectrum
as a~disjoint union of \emph{discrete} and \emph{essential spectrum},
where $\lambda\in\R$ belongs to $\sigma_{\ess}(M)$ if $\Delta-\lambda$ is not a~Fredholm operator.
Then $\lambda\in\sigma_{\d}(M)$ if and only if $\lambda$ is an eigenvalue of $\Delta$
of finite multiplicity and is an isolated point of $\sigma(M)$.
Therefore, the preimage in $L^2(M)$ of $\sigma_{\d}(M)$ under the spectral projection associated to $\Delta$
is contained in $H_{\p}(M)$.
In particular, if the essential spectrum of~$\Delta$ is empty, then $H_{\ac}(M)=H_{\sc}(M)=\{0\}$.
The essential spectrum of $\Delta$ is determined by the geometry of $M$ at infinity;
see, for example,~\cite[Proposition 3.6]{BMM17} for a general formulation of this fact.
In particular, the essential spectrum vanishes if $M$ is compact.
But there are also a~number of results about the vanishing of the essential spectrum
for non-compact manifolds;
see, for example,~\cite[Theorems 1.1 and 1.2]{BLM10} and~\cite[Examples 3.7]{BMM17}.

Now the vanishing of the essential spectrum implies the vanishing of the absolutely continuous spectrum,
which is not what we are aiming for.
In fact, we will obtain conditions which imply the absolute continuity of the spectrum, $H_{\ac}(M)=L^2(M)$,
or, in other words, the vanishing of the point and the singular continuous spectrum.
As a byproduct, we will also discuss the vanishing of the point spectrum, $H_{\p}(M)=\{0\}$.
The underlying manifolds will be \emph{Hadamard manifolds}, that is,
complete and simply connected Riemannian manifolds of non-positive sectional curvature.

